% Lecture visuelle directe avec la vision du modèle ; aucun OCR.
% Source : PDF original, vue 055, côté droit, page manuscrite 111, après le chapitre second.
% Statut : transcrit ; relecture indépendante à effectuer.
% Chapitre complet : p. 111 (fin)–127.
\chapter{Suite de la préfecture de M. Ripoud}
\begin{center}
Établissement du culte extérieur à Lyon. Premier anniversaire de la fondation.
\end{center}

\section*{Fête de la Pentecôte}
% Mention marginale de la page manuscrite 111.
Lorsque les fondateurs de la Congrégation avaient rédigé le règlement nous sortions à peine de l'état de persécution. Il n'y avait point encore de paroisses régulièrement établies ; c'est ce qui explique comment on commit la faute de comprendre des solennités comme celle de la Pentecôte parmi les fêtes que nous nous proposions de célébrer entre nous hors des assemblées des fidèles. Mais lorsqu'on vit approcher cette fête on s'aperçut bien vite que cela présentait les plus grands inconvénients ; car alors les paroisses se trouvaient parfaitement organisées et il n'eût pas été convenable dans un semblable jour de s'en isoler.

Cependant une réception était préparée ; nous en étions alors fort avides ; il eût été douloureux de la différer. Pour
% Source : PDF original, vue 056, côté gauche, page manuscrite 112.
tout concilier on convint qu'il n'y aurait pas de réunion le matin, que chacun en s'unissant en esprit à ses frères communierait dans sa paroisse et que dans la soirée à la suite des vêpres la Congrégation aurait un salut solennel. Là les récipiendaires pourraient renouveller en présence de l'adorable Eucharistie et avant de recevoir la bénédiction du St Sacrement la consécration qu'ils auraient déjà faite le matin avant la communion.

Nous dûmes à l'intervention de M. Ripoud notre préfet d'obtenir des supérieurs du séminaire la permission de nous réunir dans leur chapelle. Une grande faveur nous fut de plus accordée ; elle prouvait que la Congrégation commençait à acquérir quelque considération dans le clergé. M. Fournier, orateur fort distingué, avait cette année-là prêché le carême à St Jean. Ne se lassant pas d'évangéliser dans un moment où la parole de Dieu pouvait produire de si heureux effets il prolongea sa station jusqu'à la Fête-Dieu. Il voulut bien faire entendre sa voix puissante à notre réunion encore si obscure. Vous jugez si cette bonté de sa part nous rendit heureux. Je me souviens qu'il prononça au milieu de nous avec son éloquence ordinaire un discours sur l'établissement et la conservation de l'Église qui était bien fait pour réchauffer et soutenir notre foi.

\section*{Prédications}
% Mention marginale de la page manuscrite 112.
M. Fournier en annonçant la parole de Dieu dans notre ville au moment où nos églises venaient de se rouvrir avait rendu parmi nous un bien grand service à la cause de la religion. Le succès qu'il avait obtenu était prodigieux. Dès cinq heures du matin on se pressait autour des portes de l'église de St Jean avant même qu'elles fussent ouvertes pour se trouver
% Source : PDF original, vue 056, côté droit, page manuscrite 113.
tout prêt à envahir les places. Les personnes les plus étrangères à l'esprit religieux séduites par tout ce qu'elles en entendaient raconter accouraient pour l'entendre. Or dans ce moment où nous sortions à peine de la Révolution quoique la foi se fût grâce à Dieu conservée dans une grande partie des habitants de Lyon on ne retrouvait encore que trop souvent dans la société des traces bien vives des ravages opérés par l'esprit philosophique du dernier siècle. Il n'était pas rare de rencontrer dans les salons de ces esprits légers qui cherchaient à briller en prenant pour sujets de leurs plaisanteries également fades et impies les maximes ou les pratiques de la religion. M. Fournier prit en main avec une grande supériorité de talent la défense des vérités de la foi. Il employait si bien l'arme de l'ironie que dans la crainte de se voir écrasé sous le poids accablant du ridicule dont il les couvrait on n'aurait plus osé soutenir la cause expirante du déisme et de l'incrédulité. Sa parole était incisive et quelquefois un peu satyrique. Son éloquence était fortement sentie, la noblesse de son geste et la beauté de son organe en imposaient à l'auditoire. On peut avancer hardiment qu'il a beaucoup contribué à la révolution complète qui s'est opérée à cette époque dans l'esprit d'un grand nombre de nos salons.

Un autre prédicateur non moins distingué occupait la même année la chaire de l'église de St Louis. C'était le P. Lambert dont nous avons déjà parlé, ce bon père de la foi qui avait bien voulu nous aider de ses conseils lors de la rédaction de notre règlement. Le P. Lambert était l'orateur du sentiment. Il allait droit au cœur ; il s'exprimait avec grâce et facilité ; il s'appliquait surtout à faire apprécier à son auditoire les douceurs qu'on
% Source : PDF original, vue 057, côté gauche, page manuscrite 114.
trouve dans la pratique de la religion en même temps qu'il lui en faisait admirer les beautés. Il était préféré à M. Fournier par les âmes pieuses et tendres et d'après ce que je viens de vous en dire vous n'en serez pas étonnés. Tous ceux qui avaient conservé la foi se sentaient invinciblement entraînés par cette douce onction qui faisait tout le charme de ses discours. M. Fournier cherchait surtout à venger la religion des outrages auxquels elle avait été en butte ; par la puissance de sa parole il la plaçait sur le trône de la vérité qui n'appartient qu'à elle ; il forçait toute intelligence à se soumettre à son empire et à s'abaisser en sa présence. Le P. Lambert ouvrait à son auditoire les bras que conduisait la charité ; il plaçait cette pauvre brebis égarée sur ses épaules et à l'exemple du bon pasteur au nom duquel sa voix se faisait entendre il la portait, la soutenait, la consolait et ne l'abandonnait pas jusqu'à ce qu'enfin il eût pu la déposer dans le bercail où un charitable médecin pourra sonder la profondeur de ses blessures et y appliquer les remèdes.

Voyez, Messieurs, comme la Providence se plaisait à répandre ses bienfaits sur notre ville. Voilà deux de ses serviteurs qu'elle y fait arriver à la fois pour exercer un double ministère de miséricorde ; l'un est chargé par elle de rétablir parmi nous l'empire de la foi et de dessiller les yeux des aveugles afin qu'ils puissent enfin être éclairés par la véritable lumière ; l'autre a pour mission d'exciter les pécheurs à la pénitence et de remplir par de nouveaux adorateurs en esprit et en vérité les rangs des fidèles que les malheureux temps que nous venons de traverser n'avaient que trop éclaircis. Je
% Source : PDF original, vue 057, côté droit, page manuscrite 115.
n'ai donc pas besoin de vous ajouter que le plus que cela leur était possible les congréganistes assistaient avec empressement à ces instructions diverses. Depuis bien longtemps nous avions été privés d'entendre la parole de Dieu accompagnée de tout le charme de l'éloquence ; nous en étions donc avides car outre les avantages personnels que nous pouvions en retirer elle avait encore pour nous tout l'attrait de la nouveauté.

\section*{Rétablissement du culte extérieur}
% Mention marginale de la page manuscrite 115.
Il manquait encore quelque chose au triomphe de la religion dans notre ville. Nous ne pouvions pas franchir le seuil de nos temples. Cette privation du culte extérieur était bien cruelle pour tous les fidèles et en particulier pour les congréganistes et la présence de quelques protestants pouvait nous faire craindre d'être réduits comme Paris et plusieurs autres villes de France à ce triste état pour toujours. Nous nous réunissions bien parfois avec quelques-uns de nos frères pour aller dans les paroisses des campagnes voisines assister aux cérémonies extérieures qui du moins là pouvaient se faire en toute liberté ; mais ces jouissances passagères ne faisaient en quelque sorte qu'augmenter nos regrets. Nous voyions approcher la solennité de la Fête-Dieu. Nous étions sans doute très disposés à aller renouveller dans les cloîtres des Chartreux cette procession touchante dans laquelle la Congrégation avait pour ainsi dire pris naissance ; mais dans ce jour que nous eussions dû regarder un jour de triomphe pour notre foi nous allions encore avoir la douleur de voir nos rues et nos places privées des bénédictions que la majesté divine se plaît à répandre partout sur son passage. Oh ! que cela était pénible pour des cœurs chrétiens ! Pour des cœurs congréganistes !

% Source : PDF original, vue 058, côté gauche, page manuscrite 116.
La bonté de Dieu si souvent manifestée sur nous vint encore ici à notre secours. Aussitôt après que le cardinal Fesch eut achevé l'organisation de son diocèse il était parti pour Paris. Son neveu se disposait alors à se faire proclamer empereur. Dans des vues ultérieures dont nous verrons plus tard le développement il chargea son oncle de l'ambassade de Rome. Le cardinal dut donc se mettre en route pour se rendre auprès du pape. Il calcula sa marche de manière à arriver à Lyon quelques jours avant la Fête-Dieu et à y pouvoir séjourner jusqu'après la fête. À peine fut-il descendu de voiture qu'il fit appeler à l'archevêché le commissaire général de police ; c'était le nom qu'on donnait alors au magistrat chargé de tous les détails de la police de la ville. Dès que celui-ci se fut rendu à ses ordres : « Monsieur, lui dit-il, je vous ai fait prier de venir à l'archevêché parce que je vais vous prévenir que dimanche prochain pour la Fête-Dieu j'ai l'intention de faire la procession du St Sacrement et de plus que cette procession sera générale ; toutes les paroisses de la ville se rendront processionnellement à St Jean pour en faire partie ; ainsi prenez vos mesures en conséquence pour prévenir toute espèce de trouble. » La foudre serait tombée à côté de ce pauvre commissaire général qu'elle ne l'aurait pas plus troublé que cette déclaration du cardinal. « Monseigneur, répond-il avec émotion, votre éminence me permettra de lui observer que le culte n'est pas extérieur à Lyon et que la loi s'oppose à ce qu'il le soit parce qu'il y a des protestants. — Monsieur, vous êtes dans l'erreur ; les protestants d'une même communion ne sont point assez nombreux à Lyon pour qu'ils puissent
% Source : PDF original, vue 058, côté droit, page manuscrite 117.
légalement avoir un consistoire ; ainsi ils ne peuvent avoir aucun droit de nous empêcher de faire extérieurement les cérémonies de la religion catholique et d'y donner tout l'éclat dont elles peuvent être susceptibles. Au surplus, Monsieur, j'ai expliqué la question au premier consul et il l'a parfaitement comprise et j'ai son consentement formel. Je vous déclare donc de nouveau que la procession aura lieu ; c'est votre affaire de veiller au maintien de l'ordre. — Mais, Monseigneur, permettez-moi de vous faire remarquer que cette cérémonie va déplaire à bien du monde ; je crains que ce ne soit une occasion de désordre. Nous ne serons pas les maîtres de contenir tous ces mécontents. — Je vous répète, Monsieur, que la procession se fera. Quant aux désordres que vous craignez c'est à vous d'y pourvoir. Si le moindre trouble se manifeste tenez-vous bien pour averti, je vous en rends responsable ; je vous en préviens d'avance car votre devoir est de l'empêcher. » Le pauvre commissaire s'en fut l'oreille basse ; la fermeté du prélat l'avait déconcerté. L'oncle était tout comme le neveu : il fallait une volonté de fer. Dans sa perplexité voyant sa responsabilité mise en jeu il va consulter le préfet, le général, toutes les autorités de l'époque. S'ils avaient eu affaire à un évêque ordinaire ces Messieurs fussent peut-être bien parvenus à empêcher la procession ; mais tous furent d'avis qu'il n'était pas possible de résister au cardinal et que par conséquent il fallait bien se résigner à laisser faire. Cependant pour jeter au moins par leur conduite quelque espèce de blâme sur une manifestation qu'ils ne supportaient qu'avec la peine la plus vive ils convinrent que non seulement ils n'assisteraient pas à la procession mais que pour n'y
% Source : PDF original, vue 059, côté gauche, page manuscrite 118.
prendre part d'aucune manière ils iraient tous passer le dimanche à la campagne. Le cardinal laissa ces Messieurs bouder tout à leur aise ; seulement comme d'après un décret du premier consul il a droit en sa qualité d'ambassadeur à Rome à une garde d'honneur de trois cents hommes il la réclame pour l'escorte à l'heure fixée pour le départ de la procession et on n'osa pas la lui refuser. Il fit connaître aux fidèles par l'intermédiaire de MM. les curés les rues et les places que la procession devait parcourir et les fit inviter à tapisser la façade de leurs maisons. De plus il ordonna au clergé des paroisses de se rendre processionnellement avec leurs paroissiens à la cathédrale au jour et à l'heure indiquée pour la cérémonie.

Toute la ville attendait avec curiosité l'issue de cette lutte ouverte entre la religion et les derniers débris de ces hommes qui avaient mis tant de persistance à faire cesser parmi nous les fêtes du Seigneur. Nous étions pleins de confiance ; la haute protection du cardinal et bien plus que cela la bonté de Dieu qui tant de fois déjà s'était manifestée en notre faveur devait nous faire espérer la victoire. Nous ne pouvions pas cependant être tout à fait sans crainte. C'est donc ballottés entre ces deux sentimens que nous vîmes enfin arriver le beau jour de la Fête-Dieu. M. Fournier devait ce jour-là prêcher pour la dernière fois à St Jean après la grand-messe. Comme à l'ordinaire son auditoire était immense. Nos congréganistes s'y trouvaient en grand nombre ; l'orateur monte en chaire, se recueille un instant, promène un regard assuré sur cette imposante assemblée, fait le signe de la croix et d'une voix ferme et majestueuse fait entendre ces paroles du psalmiste qu'il
% Source : PDF original, vue 059, côté droit, page manuscrite 119.
a choisies pour son texte : \textit{Exurgat Deus et dissipentur inimici ejus et fugiant omnes qui oderunt eum a facie ejus}\footnote{Traduction littérale : « Que Dieu se lève, que ses ennemis soient dispersés et que tous ceux qui le haïssent fuient devant sa face. » Psaume 67 (68), verset 2. (Note du transcripteur).}. « Que le Seigneur se lève, que ses ennemis soient confondus et que tous ceux qui le haïssent fuient devant lui. » Ces paroles dans la circonstance étaient par elles-mêmes assez significatives. La noble attitude du prédicateur, le ton dont il les prononça, le geste expressif dont il les accompagna firent assez comprendre quels étaient ces ennemis que la présence adorable du Dieu tout-puissant allait confondre et qui se trouveraient réduits pour l'éviter à prendre la fuite. L'orateur compare ensuite dans un exorde pompeux le triomphe de l'arche sainte revenant victorieuse de sa captivité chez les Philistins sur lesquels dans sa vengeance la main du Seigneur s'était appesantie pour les forcer à la renvoyer à son peuple et cet autre triomphe solennel que la population de notre ville préparait en ce moment à J.-C. dans le sacrement de son amour. À cet exorde qui fit palpiter tous nos cœurs succéda un sermon rempli d'éloquence sur l'Eucharistie. Pleins de l'émotion que ce sermon avait excitée en nous dès qu'il est terminé nous nous pressons de sortir de l'église pour aller nous disposer à assister à la grande et belle cérémonie qui doit commencer dans quelques heures.

Quel spectacle admirable nous présente l'aspect de la ville en ce moment ! Aucun ordre de police, pas même la moindre invitation n'était intervenue de la part de l'autorité et de tous les côtés on voyait les habitants empressés à balayer les rues, à tapisser les maisons, à élever des reposoirs, à joncher le pavé de fleurs ou à tresser des guirlandes. Une pensée unique paraissait animer toute la population. La plus douce sérénité brillait sur tous les
% Source : PDF original, vue 060, côté gauche, page manuscrite 120.
visages. Seuls au milieu de ce joyeux concours la physionomie rembrunie, l'œil sauvage, nos magistrats comme s'ils eussent voulu se charger eux-mêmes de vérifier le texte trop bien appliqué du prédicateur montaient en voiture et fuyaient...

\section*{Procession de la Fête-Dieu}
% Mention marginale de la page manuscrite 120.
Cependant de tous les points de la cité les paroisses se sont mises en mouvement et se dirigent sur St Jean. Après une interruption si longue et si cruelle la croix de J.-C. reparaît enfin dans les rues de notre ville triomphante et suivie par de nombreux chrétiens dont elle est le drapeau. À deux heures tous les fidèles sont réunis dans la métropole et la procession s'ébranle. Le St Sacrement est porté par le cardinal au milieu d'un nombreux clergé ; il est précédé et suivi par de longues lignes de fidèles des flambeaux à la main. Non seulement la Congrégation était là tout entière mais elle avait usé de son mieux de la petite influence qu'elle commençait à acquérir pour augmenter le nombre des enfants de la foi qui venaient former le cortège de notre divin Sauveur. Parmi tous les hommes en place un seul, l'adjoint de Vaise, marchait en costume à la suite du St Sacrement. Rendons hommage à ce digne magistrat dont je regrette de ne pas savoir le nom pour le citer ici, qui supérieur à tout respect humain n'a pas craint dans cette circonstance de manifester ouvertement sa foi.

La procession parcourut pendant six heures entières une grande partie de la ville. Partout une foule immense se prosternait sur le passage de l'auguste sacrement ; partout l'allégresse était générale. Si çà et là quelques impies de bas étage essayaient d'insulter à la joie
% Source : PDF original, vue 060, côté droit, page manuscrite 121.
commune et de braver la foi des fidèles en restant insolemment couverts au moment où ils se trouvaient en présence du corps adorable de J.-C., ils étaient si peu nombreux qu'ils pouvaient bien passer inaperçus ; le peuple en était irrité mais la vengeance qu'il exerçait sur eux était bien innocente. Il poussait ces malheureux dans les allées des maisons, fermait sur eux les portes pour mettre fin au trouble qu'ils occasionnaient et leur adressait en même temps ces paroles : « Vous vous croyez encore en 93 ; vous vous trompez, grâce à Dieu le temps de sinistre mémoire est passé. »

Dès que la procession fut rentrée Monseigneur réunit MM. les curés à l'archevêché et leur adressa ces paroles : « Messieurs, chacun de vous fera dimanche prochain la procession dans sa paroisse. Mettez-vous en mesure pour être pourvus d'ici là de tout ce qui est nécessaire. Donnez à cette cérémonie autant d'éclat que cela vous sera possible. Désormais si il se présente des enterrements vous les ferez processionnellement dans les rues ; désormais aussi vous porterez publiquement le bon Dieu aux malades. Vous m'objecterez peut-être que vous n'avez pas de dais ; ce ne serait pas une raison pour retarder, dans ce cas jusqu'à ce que vous ayez pu vous en procurer un vous porterez le St Sacrement sous un parapluie car il importe que vous vous mettiez promptement en possession pendant que je suis encore ici pour vous défendre. »

Tout s'exécuta à la lettre suivant les ordres du cardinal et personne n'osa y faire aucune opposition. Les protestants eux-mêmes n'essayèrent pas de se plaindre. Ainsi le dimanche qui terminait l'octave de la Fête-Dieu renouvela le spectacle
% Source : PDF original, vue 061, côté gauche, page manuscrite 122.
que la fête elle-même avait offert. L'élan qu'il occasionna devint même peut-être plus général encore. Car les processions se trouvant aussi nombreuses que les paroisses de la ville il en résulta que les rues tapissées qu'elles devaient parcourir se croisaient dans toutes les directions et que les reposoirs multipliés à l'infini s'élevaient sur toutes les places. Le zèle de la population n'en parut que plus vif. Du reste, Messieurs, je ne vous raconte ici que ce que vous avez vu naguères de vos propres yeux lorsque les regrets vivement exprimés des habitants de Lyon ont forcé à les laisser enfin jouir paisiblement de cette solennité qu'on avait encore une fois essayé de leur ravir.

\section*{Anniversaire de la fondation}
% Mention marginale de la page manuscrite 122.
Ce fut peu de jours après ce jour de triomphe qu'un anniversaire bien précieux pour nous vint réveiller et exciter de plus en plus dans nos cœurs les sentimens déjà si vifs qui nous attachaient par des liens indissolubles à notre chère Congrégation. Vous me prévenez, Messieurs, et vous comprenez déjà que je vais vous parler de la première célébration solennelle du troisième dimanche de juillet. Si l'annonce de chaque nouvelle fête de la Congrégation nous remplissait par avance d'une joie bien pure vous devez juger quel effet devait produire celle d'une solennité où en fêtant pour la première fois le St Cœur de Marie nous allions donner un libre cours à l'effusion de notre reconnaissance envers notre bon Sauveur qui pour assurer notre salut par la voie de l'amour avait voulu nous réunir tous dans le cœur de sa mère.

Un article du règlement prescrivait de célébrer ce
% Source : PDF original, vue 061, côté droit, page manuscrite 123.
grand jour par une réunion qui devait durer toute la journée. Nous étions bien résolus d'exécuter cet article rigoureusement à la lettre ; la seule difficulté était de trouver un local convenable. MM. Paul et de Villers vinrent à notre secours. M. Paul, vénérable et saint ecclésiastique, était curé des Chartreux. M. de Villers ancien vicaire général y avait fixé sa résidence et y exerçait le saint ministère. Ces MM. voulurent bien nous offrir de nous recevoir ; ils nous mirent pour la première fois en possession de cette petite chapelle que les congréganistes ont fait retentir si souvent depuis des louanges de Dieu et où tant de grâces ont découlé avec abondance sur leurs cœurs.

Dès le grand matin la Congrégation fut réunie ; elle se trouva composée de 32 membres présents. On commença comme à présent par la récitation de l'office puis le \textit{Veni Creator}\footnote{Début de l'hymne latine : « Venez, Créateur. » (Note du transcripteur).}, ensuite la grand-messe. Les fondateurs et les nouveaux reçus présentèrent ensemble le pain bénit : les premiers et les derniers. Cela renfermait toute la Congrégation ; tous comme ils en ont conservé jusqu'à présent l'heureuse habitude à toutes leurs fêtes les congréganistes s'unirent de nouveau de la manière la plus intime en participant tous ensemble au sacrement adorable de nos autels. Oh ! Messieurs, comme la religion aggrandit l'esprit de l'homme ! Que de charmes indicibles elle ajoute à tous nos sentimens ! Quelle vie que la vie de foi ! Elle donne à l'âme une vigueur nouvelle ; elle la rend plus capable d'aimer et de sentir. Elle retrempe les liens de l'amitié dans cette charité qui fait son essence ; elle fait éprouver à ceux qui la possèdent des jouissances telles que ceux qui y sont étrangers ne peuvent
% Source : PDF original, vue 062, côté gauche, page manuscrite 124.
pas même en concevoir la pensée. Par elle les amis dont elle a fait des frères sont mis en possession de ce bonheur ineffable que le St Esprit a promis aux frères qu'une même pensée d'amour de Dieu fait véritablement habiter ensemble et surtout lorsque par la communion le cœur de notre divin maître devient au même moment la demeure de tous. Oh ! c'est bien alors que nous pouvons nous écrier avec transport : \textit{Ecce quam bonum et quam jucundum habitare fratres in unum}\footnote{Traduction littérale : « Voici combien il est bon et agréable que des frères habitent ensemble. » Psaume 132 (133), verset 1. (Note du transcripteur).}. Quelle union en effet que celle que Dieu vient ainsi former lui-même et dont il veut être à la fois le principe et le centre ! Toutes ces réflexions occupaient délicieusement les cœurs des congréganistes au moment où ils s'approchèrent du saint autel pour rendre grâces au Seigneur de tous les bienfaits déjà répandus sur eux en s'unissant à la Ste Eucharistie seul sacrifice d'actions de grâces qui puisse être digne d'être offert à la majesté de notre Dieu.

Après la grand-messe eurent lieu les cérémonies en usage dans les réceptions et la bénédiction du St Sacrement.

On passa ensuite dans la salle du déjeuner. Là le pain bénit de 1802 vint se placer à côté de celui de 1803 ; l'un et l'autre furent distribués ensemble à tous les frères.

Le reste de la matinée fut employé en conversations amicales ; elles furent cependant interrompues pour quelques instants par une assemblée générale qu'on voulut tenir en mémoire de la première qui avait eu lieu le jour même de la fondation. On y lut le procès-verbal de cette première journée de notre existence. Le préfet y proclama ensuite la liste des candidats
% Source : PDF original, vue 062, côté droit, page manuscrite 125.
pour la nomination de son successeur, usage qui s'est perpétué jusqu'à ce jour dans la Congrégation.

Vers midi la Congrégation se réunit autour d'une table commune. Le P. Roger et le préfet s'étaient chargés de faire les honneurs. MM. Paul et de Villers avaient bien voulu prendre place à ce banquet fraternel ; il est difficile même pour ceux qui ont assisté plusieurs fois à cet anniversaire de la fondation de se figurer tout ce qu'éprouvaient nos cœurs en ce beau jour. Il était si consolant pour nous d'avoir déjà à la fin d'une première année toutes nos espérances remplies. Non seulement nous avions rencontré dans l'établissement de la Congrégation les avantages dont nous nous étions flattés mais elle était pour nous une source de jouissances sans cesse renouvellées auxquelles nous étions bien loin de nous attendre. Aussi comme nous en savourions les douceurs ! Nos âmes se livraient en toute liberté à des sentimens de bonheur d'autant plus précieux qu'ils étaient sans remords. Le P. Roger qui savait mieux que personne le secret d'exciter en nous cette aimable gaîté qui convient si bien aux serviteurs de notre bon maître ne négligeait aucun moyen pour la rendre plus vive. Vous vous intéressez trop à ce qui concerne la Congrégation pour que je ne me fasse pas un devoir d'entrer avec vous jusque dans les plus petits détails. Je ne craindrai donc pas de vous dire que non seulement on porta des toasts d'abord au souverain pontife puis à nos vénérables hôtes, à notre père, au préfet, aux fondateurs et enfin à la Congrégation toute entière, mais encore le P. Roger nous apprit une espèce de marche au moyen de laquelle chacun venait en signe d'union trinquer avec tous
% Source : PDF original, vue 063, côté gauche, page manuscrite 126.
ses frères. Cette cérémonie s'est conservée dans la Congrégation jusqu'au moment où les réunions sont devenues si nombreuses qu'elles en ont rendu la pratique impossible.

Après le repas on rentra dans la chapelle pour chanter solennellement les vêpres et les litanies de la Ste Vierge, renouveller tous ensemble notre consécration à Marie et terminer le tout par le \textit{Te Deum}\footnote{Début de l'hymne latine : « Vous, Dieu, [nous vous louons] ». (Note du transcripteur).} suivi de la bénédiction du St Sacrement. Après avoir achevé ainsi la partie pieuse de la fête on employa la soirée à se livrer avec simplicité à toute l'effusion d'une joie de famille. Mais auparavant on se souvint du baiser de charité que s'étaient mutuellement donné les fondateurs au jour de la fondation. On établit l'usage de renouveller chaque année entre tous les frères cette expression des sentimens qui les unissent. Notre père et nos hôtes furent ensuite les premiers à nous donner l'exemple de nous abandonner librement à tout le charme des plaisirs de l'amitié. M. de Villers surtout paraissait jeune de se trouver au milieu de cette jeunesse ; il se prêtait à toutes nos folies avec l'ardeur d'un jeune homme. Les chansons vinrent comme cela s'est toujours pratiqué depuis embellir encore cette belle journée. M. Chaulet qui bientôt va jouer un rôle important dans cette histoire fut notre premier poète. Les anciens membres n'ont pas pu oublier comme je ne l'ai pas oublié moi-même ce refrain que nous avons répété si longtemps avec tant d'enthousiasme :
\begin{quote}
De notre auguste patronne\par
Imitons l'amour ardent,\par
Pour obtenir la couronne\par
Allons toujours en avant.
\end{quote}
Le mot \textit{en avant} devint comme notre cri de guerre. C'était avec ce Montjoie St Denys
% Source : PDF original, vue 063, côté droit, page manuscrite 127.
de la Congrégation que le P. Roger nous excitait à ne jamais rester en arrière dans le service de Dieu. En avant, répétait-il souvent, en avant jusqu'au ciel !

Rappelons en passant puisque nous avons eu occasion de parler des travaux poétiques de M. Chaulet que c'est à lui que la Congrégation a l'obligation de ce cantique des réceptions que vous chantez encore avec tant de plaisir et qui dans son aimable simplicité produit toujours tant d'effet soit sur l'enfant de Marie qui vient prendre sa place dans la famille soit sur tous les frères qui s'apprêtent à le serrer contre leur cœur dans toute l'effusion de la charité.

La célébration de ce premier anniversaire eut pour la Congrégation des résultats si avantageux qu'ils dépassèrent toutes nos prévisions ; elle sembla acquérir une nouvelle vie ; les lâches s'armèrent d'un courage tout nouveau, les tièdes furent réchauffés, tous en général retrouvèrent leur vigueur première. Nous aurons bientôt à en admirer les effets. Conformément au règlement l'élection du préfet eut lieu le jour de la fête de l'Assomption. Les études ecclésiastiques de M. Ripoud ne permettaient pas de le conserver plus longtemps à la tête de la Congrégation. M. Chaulet fut appelé par les voix de ses confrères à lui succéder. La Congrégation voulut néanmoins donner à M. Ripoud un témoignage de sa reconnaissance pour tous les services qu'il lui avait rendus pendant cette première année. Par un arrêté \textit{ad hoc} elle lui conféra le titre honorifique de premier préfet de la Congrégation\footnote{\textit{Ad hoc} : « pour cela », ici « pris spécialement à cet effet ». (Note du transcripteur).}.
